% ======================================================================
% Chapter 3 -- Part C: money in the trade-off table, rollout, memo,
% end matter
% ======================================================================

\section{Can We Put Dollars on the Trade-off?}
\label{sec:money-tradeoff}

Chapter~\ref{ch:toolkit} scored designs on a 1--5 scale for all eight criteria\index{evaluation criteria}. With the tools of this chapter, three of the criteria can usually be replaced by money per transaction: \textbf{cost} becomes the run-cost term of Equation~\eqref{eq:unit}, \textbf{human workload} becomes the human-time term, and much of \textbf{quality} becomes the expected-error-cost term. The other five---speed, reliability, security and privacy, control and auditability, maintainability---usually stay as scores, because their value is either already inside a KPI (speed drives Harbor's win rate) or shows up mainly as risk.

A mixed table, with a dollar row and score rows, is more honest than an all-score table. It shows the sponsor exactly what money is at stake, and keeps the judgment calls visible as judgment calls. Table~\ref{tab:money} applies the idea to Team~A's and Team~B's quote designs from Chapter~\ref{ch:foundations}.

\begingroup
\small
\begin{longtable}{@{}lrr@{}}
\caption{Team A's and Team B's quote designs with money per quote. Error rates for Team A are the pilot estimates from a 200-quote test; human-time figures from the analysts' time log. Fixed platform and support costs, about \$4.20 per quote for either design, are left out because they do not differ.}\label{tab:money}\\
\toprule
\rowcolor{coHead} Per quote & \cellcolor{coRedBg}Team A (single agent) & \cellcolor{coGreenBg}Team B (workflow) \\
\midrule
\endfirsthead
\toprule
\rowcolor{coHead} Per quote & \cellcolor{coRedBg}Team A (single agent) & \cellcolor{coGreenBg}Team B (workflow) \\
\midrule
\endhead
Model and tool calls & \$0.21 & \$0.04 \\
Analyst and manager time & \$12.10 & \$10.50 \\
Expected error cost, below-floor quotes & \$4.28 & \$0.71 \\
Expected error cost, wrong contract & \$1.90 & \$0.60 \\
\midrule
\rowcolor{coAmberBg} Total cost per quote & \$18.49 & \textbf{\$11.85} \\
\midrule
Speed (score) & 5 & 4 \\
Control and auditability (score) & 1 & 5 \\
Security and privacy (score) & 2 & 4 \\
\bottomrule
\end{longtable}
\endgroup

Team A's design is more than five times as expensive to run, but that is not why it loses: run cost is still a few cents next to the human time and error cost. It loses because its error cost is several times higher---a policy in a prompt is broken more often than a policy in code---and because no one can audit why it priced a quote as it did. The dollar rows make the first point; the score rows make the second.

% ----------------------------------------------------------------------
\section{How Do We Try It Safely?}
\label{sec:rollout}

\begin{thinkabout}
Imagine the pilot ends and the numbers look wonderful. What could make those wonderful numbers misleading? List two possibilities before you read the six parts of a good pilot plan.
\end{thinkabout}

No sponsor should be asked to approve a full launch on the strength of a design document. The usual request is for a \keyterm{pilot}\index{pilot|textbf}: a limited, time-boxed trial with a decision at the end. A good pilot plan has six parts.

\begin{enumerate}[itemsep=2pt]
  \item \emph{Scope.} Who and what is in the pilot, and who is not. Harbor: two of six sales regions; quotes up to \$250{,}000; standard catalog items only.
  \item \emph{Comparison.} What the pilot results will be compared against. Harbor: the four non-pilot regions over the same weeks, and the same two regions over the previous year's same weeks. Without a comparison group, a seasonal upswing will be credited to the AI.
  \item \emph{Duration.} Long enough to measure the most uncertain benefit. Win rate on large quotes moves slowly; Harbor's analysts estimated that eight weeks would give roughly 1{,}000 pilot quotes, enough to detect a two-point change in win rate with reasonable confidence but probably not a one-point change. The plan says so: the pilot can confirm the floor case, and can only bound the growth case.
  \item \emph{Go/no-go criteria.}\index{go/no-go criteria|textbf} Stated in advance, measurable, and tied to KPIs and error budgets (Table~\ref{tab:gonogo}). Criteria agreed after the results are known are not criteria.
  \item \emph{Stop rules.}\index{stop rule|textbf} Conditions that halt the pilot at once, before the scheduled decision: any below-floor quote issued; any customer data exposure.
  \item \emph{Owner.} One named person who makes the go/no-go call and owns the solution afterward. Committees review; owners decide.
\end{enumerate}

\begingroup
\small
\begin{xltabular}{\textwidth}{@{}>{\raggedright\arraybackslash}X>{\raggedright\arraybackslash}p{2.6cm}>{\raggedright\arraybackslash}p{3.2cm}@{}}
\caption{Go/no-go criteria for Harbor's eight-week pilot.}\label{tab:gonogo}\\
\toprule
\rowcolor{coHead} Criterion & Baseline & Go if \\
\midrule
\endfirsthead
\toprule
\rowcolor{coHead} Criterion & Baseline & Go if \\
\midrule
\endhead
Median quote turnaround & 26 h & $\le$ 6 h (target 4 h) \\
Analyst minutes per quote & 35 min & $\le$ 12 min \\
Realized margin, pilot vs.\ comparison regions & --- & No worse than $-0.3$ points \\
Below-floor quotes issued & 2.4\% & 0 (stop rule) \\
Wrong customer or contract & 0.9\% & $\le$ 0.3\% \\
Share of eligible quotes sent through the system & --- & $\ge$ 70\% by week 6 \\
Representative satisfaction (weekly 1--5 pulse) & --- & $\ge$ 3.8 \\
\bottomrule
\end{xltabular}
\endgroup

\begin{watchout}[Watch out: the easy-quotes pilot]
Note the adoption criterion in Table~\ref{tab:gonogo}. A pilot in which representatives route only the easy quotes through the system will show excellent numbers that will not survive a full launch.
\end{watchout}

After the pilot, the decision has three outcomes, not two: \cotag{coGreen}{Go} means expand on a stated schedule, \cotag{coAmber}{Iterate} means fix named problems and repeat a shorter pilot, and \cotag{coRed}{No-go} means stop and record why. Planning for all three in advance keeps the pilot from drifting into a permanent half-launched state, which is the most common fate of AI pilots that have no owner.

% ----------------------------------------------------------------------
\section{How Do We Ask for a Yes?}
\label{sec:memo}

Every case in this book ends with a memo to the executive sponsor, and every memo follows the same outline.\index{one-page memo|textbf} The limit of one page is a discipline, not a formatting rule: if the case cannot be made in a page, it is not yet clear, and an executive will not read the second page anyway. Figure~\ref{fig:memo} is Harbor's.

\begin{figure}[htbp]
\begin{tcolorbox}[colback=white, colframe=black!60, boxrule=0.6pt, arc=1pt, left=10pt, right=10pt, top=8pt, bottom=8pt, fontupper=\small]
\textbf{To:} Vice President, Sales \hfill \textbf{From:} Quote Design Team\\
\textbf{Re:} Faster quotes with margin protection --- request to pilot \hfill \textbf{Decision needed by:} 14 March

\smallskip\hrule\smallskip

\textbf{Recommendation.} Approve an eight-week pilot of the redesigned quote process in the Central and Gulf regions, at a build cost of \$180{,}000.

\textbf{The problem.} Large quotes take a median of 26 hours, almost all of it waiting in queues. 2.4\% of quotes are issued below the margin floor. Representatives report losing deals to faster competitors.

\textbf{What changes.} Requests are read and structured automatically; prices, contract terms and policy checks are applied by tested rules; analysts check drafts instead of building quotes; managers see only true exceptions, with the reason stated. Representatives still send every quote.

\textbf{Value.} \emph{Floor case, from measured inputs:} about \$335{,}000 a year (analyst capacity and fewer below-floor quotes, net of run costs); payback in about seven months. \emph{Upside, to be tested:} each one-point rise in win rate on quoted orders adds about \$2.3 million in gross profit a year.

\textbf{Risks and safeguards.} Below-floor quotes are blocked in code; any occurrence stops the pilot. Customer data stays within our cloud tenancy. Adoption is the main risk; the pilot measures it directly.

\textbf{Pilot and decision.} Comparison against the four other regions. Go if median turnaround $\le$ 6 hours, analyst time $\le$ 12 minutes per quote, margin within 0.3 points of comparison regions, wrong-contract errors $\le$ 0.3\%, and at least 70\% of eligible quotes use the system. Decision owner: VP Sales. Expansion to all regions would follow in one quarter.

\textbf{What we need from you.} Approval of the budget; two regional managers as pilot leads; a decision on whether analysts' freed time goes to the strategic-accounts team.
\end{tcolorbox}
\caption{Harbor's one-page memo to the executive sponsor. The template is in Appendix~C.}
\label{fig:memo}
\end{figure}

\begin{tryit}
Read the memo in Figure~\ref{fig:memo} as the vice president would, in under two minutes. Can you say what she is being asked to approve, what it costs, what she gets, and what could go wrong? If you can, the memo works.
\end{tryit}

Five features of the memo are deliberate and should be copied:
\begin{itemize}[itemsep=1pt]
  \item The \emph{recommendation comes first}, with the money and the decision requested.
  \item The \emph{problem is stated in the sponsor's numbers}, from the process map.
  \item \emph{What changes} is described in business terms. Words like \textbf{agent}, \textbf{RAG\index{retrieval-augmented generation (RAG)}} and \textbf{MCP\index{Model Context Protocol (MCP)}} do not appear; the design table is an attachment for those who want it.
  \item The \emph{value is split into floor and upside}, with the source of each.
  \item The memo \emph{ends with what the sponsor must do}, including the decision that only she can make: what to do with the freed capacity, without which the labor benefit stays on paper.
\end{itemize}

% ----------------------------------------------------------------------
\section{The Business Layer in Every Case}
\label{sec:business-in-cases}

From Chapter~\ref{ch:case-policy} onward, each case applies this chapter in the same places of the chapter template:
\begin{itemize}[itemsep=1pt]
  \item Section 1, \emph{business brief and current process}: the process map with times, costs and error points.
  \item Section 2, \emph{stakeholders, KPIs and value hypothesis}: the four stakeholder roles, a KPI tree, and the hypothesis in one sentence.
  \item Section 4, \emph{trade-offs}: unit economics per design, with money rows where they can be measured.
  \item Section 8, \emph{business case and rollout}: the value case with its tornado analysis, error budgets, the pilot plan and the memo.
\end{itemize}
Each case also has a \textbf{business skill} it emphasizes. Chapter~\ref{ch:case-policy} focuses on defining KPIs for a solution whose output is information rather than transactions; Chapter~\ref{ch:case-support} on unit economics at high volume; and so on through the book.

% ----------------------------------------------------------------------
\section*{Chapter Summary}
\addcontentsline{toc}{section}{Chapter Summary}

\begin{itemize}[itemsep=2pt]
  \item A current-process map with one row per step and its touch times, wait times and error points is the baseline for every claim a design makes. Measure it before the pilot; it cannot be measured afterward.
  \item Sponsors, users, approvers and owners each need different things; each approver appears in the design table as a person who decides or as a constraint, and each owner is named next to what they own.
  \item Good KPIs are business-facing, measured today, movable by the design, paired with a quality guard, and few. A KPI tree links them to the metrics of individual steps in the design.
  \item A value case separates growth, labor capacity, error reduction and risk reduction. Lead with the floor case built from measured inputs; use tornado analysis to find the assumption the case rests on, and design the pilot to test it.
  \item Unit economics (Equation~\ref{eq:unit}) give a cost per transaction by design and by path, including expected error cost. Model cost is often the smallest term.
  \item When error costs are asymmetric, choose thresholds that minimize expected cost, not error count; where a rule can decide exactly, prefer the rule. Error budgets turn quality hopes into commitments.
  \item A pilot needs scope, a comparison group, enough duration, go/no-go criteria agreed in advance, stop rules and a single owner. The one-page memo makes the case to the sponsor in the sponsor's terms.
\end{itemize}

\section*{Key Terms}
\addcontentsline{toc}{section}{Key Terms}

\begin{multicols}{2}
\raggedright\small
asymmetric error cost \\
current-process map \\
error budget \\
floor case \\
go/no-go criteria \\
KPI tree \\
key performance indicator (KPI) \\
one-page memo \\
pilot \\
risk register \\
sponsor, users, approvers, owners \\
stop rule \\
tornado analysis \\
touch time and wait time \\
unit economics \\
value case \\
value hypothesis
\end{multicols}

\section*{Review Questions}
\addcontentsline{toc}{section}{Review Questions}

\begin{enumerate}[itemsep=2pt]
  \item Why does Harbor's process map change the designer's priorities, even though the analysts' work takes only 35 minutes?
  \item Name the four stakeholder roles in Harbor's case. Where does each of them appear in the design table?
  \item Give an example of an unpaired KPI for an AI solution and explain how the solution could improve it while making the business worse.
  \item What is the difference between labor capacity and cash savings? What decision turns one into the other?
  \item In Harbor's tornado analysis, why does model cost barely matter? Describe a business where it would dominate.
  \item Explain Equation~\eqref{eq:threshold} in words. Why should the threshold for sending quotes to a manager differ from the accuracy-maximizing threshold?
  \item What is an error budget, and how does one become a go/no-go criterion?
  \item Why does the pilot need a comparison group? What could mislead Harbor without one?
  \item Why does the memo avoid words like \textbf{agent} and \textbf{MCP}?
  \item A pilot ends with mixed results: turnaround and analyst time pass, adoption reaches only 55\%. What should the owner decide, and what information would help?
\end{enumerate}

\section*{Exercises}
\addcontentsline{toc}{section}{Exercises}

\paragraph{Design track}
\begin{enumerate}[itemsep=2pt]
  \item \textbf{Map a process.} For the process you have used in earlier exercises, build a current-process map with touch times, wait times and at least two error points. State how you would measure each number.
  \item \textbf{KPI tree.} Build a KPI tree for Harbor's returns process from Chapter~\ref{ch:toolkit}, with technical metrics that belong to steps of the straight-through design.
  \item \textbf{Price the errors.} For the returns case, estimate $L_k$ for an out-of-warranty credit and for a valid return wrongly refused. Which is more expensive, and what does that imply for the safe-zone rules?
  \item \textbf{Write the memo.} Write the one-page memo for Harbor's straight-through returns design, with a floor case and an upside.
\end{enumerate}

\paragraph{Build track}
\begin{enumerate}[resume, itemsep=2pt]
  \item \textbf{Probabilistic value case.} Replace the one-at-a-time sensitivity in Listing~\ref{lst:value} with a Monte Carlo simulation in which each uncertain input is drawn from a distribution you justify. Report the probability that net annual value is below \$250{,}000.
  \item \textbf{Threshold under asymmetric costs.} Using the synthetic routing scores in the repository, plot $E(\tau)$ from Equation~\eqref{eq:threshold} for $L_{\text{FN}}/L_{\text{FP}} \in \{1, 10, 100\}$ and report the optimal threshold for each ratio.
\end{enumerate}

% ----------------------------------------------------------------------
\section*{Lab 3: A Value Case You Can Defend}
\addcontentsline{toc}{section}{Lab 3: A Value Case You Can Defend}

\paragraph{Goal.} Produce a value case, a pilot plan and a memo for Harbor's returns design from Chapter~\ref{ch:toolkit}, using data rather than guesses wherever data exists.

\paragraph{Time.} About 90 minutes.

\paragraph{Steps.}
\begin{enumerate}[itemsep=2pt]
  \item Open \texttt{ch03/} in the companion repository. It contains a synthetic extract of 12 months of Harbor's returns: case timestamps, representative handling times, credit amounts, and an audit sample of 400 credits marked in or out of warranty.
  \item Compute the baseline: median and 90th-percentile days to RMA and to credit, representative minutes per return, and the out-of-warranty credit rate with a 95\% confidence interval.
  \item Copy \texttt{value\_case.py} to \texttt{returns\_value.py} and rewrite the inputs for the returns case. Mark each input \texttt{MEASURED} or \texttt{ASSUMED} in a comment.
  \item Run the tornado analysis. Identify the assumption the case rests on.
  \item Write go/no-go criteria for a six-week pilot, including at least one error budget and one adoption criterion.
  \item Write the one-page memo.
\end{enumerate}

\paragraph{Deliverable.} The baseline table, the value case with its tornado output, the go/no-go table and the memo. In one paragraph, explain how the pilot would measure the assumption your case rests on.

% ----------------------------------------------------------------------
\section*{Further Reading}
\addcontentsline{toc}{section}{Further Reading}

\begin{itemize}[itemsep=2pt]
  \item \citet{kohavi2020trustworthy} is the standard practical guide to controlled experiments; read its chapters on choosing metrics and on guardrail metrics before designing a pilot comparison.
  \item \citet{provost2013datascience} explains expected-value thinking for decisions under uncertainty, including cost-sensitive thresholds, in business language.
  \item \citet{parmenter2019kpi} covers the selection and use of KPIs in organizations, with many examples of KPIs that backfired.
\end{itemize}
